\documentclass[10pt,twocolumn]{article}
\usepackage[margin=0.72in, top=0.65in, bottom=0.65in]{geometry}
\usepackage{amsmath, amssymb}
\usepackage{booktabs}
\usepackage{xcolor}
\usepackage{titlesec}
\usepackage{enumitem}
\usepackage{microtype}
\usepackage{parskip}
\usepackage{mdframed}
\usepackage{fancyhdr}
\usepackage{hyperref}

\definecolor{accent}{RGB}{31,78,121}
\definecolor{accentblue}{RGB}{46,117,182}
\definecolor{lightblue}{RGB}{214,228,240}

\titleformat{\section}{\normalsize\bfseries\color{accent}}{}{0em}{}[\vspace{-4pt}\textcolor{accentblue}{\rule{\linewidth}{0.5pt}}]
\titleformat{\subsection}{\small\bfseries\color{accentblue}}{}{0em}{}
\titlespacing{\section}{0pt}{11pt}{4pt}
\titlespacing{\subsection}{0pt}{8pt}{3pt}

\hypersetup{colorlinks=true,linkcolor=accentblue,urlcolor=accentblue}

\pagestyle{fancy}
\fancyhf{}
\rhead{\scriptsize\color{gray}CS109 Challenge Project --- Ajay Shah}
\rfoot{\scriptsize\color{gray}\thepage}
\renewcommand{\headrulewidth}{0.3pt}

\newmdenv[
  backgroundcolor=lightblue!50,
  linecolor=accentblue,
  linewidth=0.6pt,
  innerleftmargin=7pt, innerrightmargin=7pt,
  innertopmargin=6pt,  innerbottommargin=6pt,
  skipabove=5pt, skipbelow=5pt
]{infobox}

\setlength{\parskip}{4pt}
\setlength{\columnsep}{18pt}
\setlist[itemize]{itemsep=2pt, topsep=3pt, leftmargin=14pt}

\begin{document}

\twocolumn[{%
\begin{center}
  {\large\bfseries\color{accent}F1 Race Intelligence: A Live Probabilistic Engine}\\[4pt]
  {\small\color{accentblue}CS109 Challenge Project\quad|\quad Ajay Shah\quad|\quad Stanford University\quad|\quad 2026}
\end{center}
\vspace{2pt}\noindent\rule{\linewidth}{0.8pt}\vspace{6pt}
}]

\section{Abstract}

This project is a live probabilistic system that estimates each driver's win
probability in a Formula~1 race, updating after every lap. It takes in actual race data
from the OpenF1 API and does a three-stage probability analysis: Gaussian MLE for
corner execution scoring, Naive Bayes lap classification, and a
Bayesian updater that fuses four evidence sources. It then computes,
for each driver at each lap, which single move would increase their win probability the most. The system is
designed as a proof-of-concept for Stanford's solar car racing team, where the
same framework with telemetry variables swapped provides real-time strategic
guidance during competition.

\section{Problem Statement \& Motivation}

A Formula~1 race produces $\approx$1,100 laps of telemetry across 22 drivers,
and the API samples the following data every 3 seconds: speed, throttle, brake, GPS, gap to leader, and race
position. The question I want to answer is, given all observations in a lap $k$, what is
$P(\text{winner} = i)$ for each driver $i$?

This felt like a natural Bayesian inference problem to me. I start with a prior over
outcomes, observe sequential evidence each lap, and maintain a posterior
distribution. The challenge was combining four evidence sources:
position and gap are strong late-race signals, while lap
time is more informative early when positions are still unsettled due to pit
strategy and opening-lap incidents. After talking with Juliette about my project,
I realized I could also add corner execution to my data, increasing accuracy but also
the complexity of the model.

My motivation for building this is Stanford's solar car team. Solar racing involves identical
strategic tradeoffs: energy conservation vs.\ speed, segment-by-segment
benchmarking, and real-time decisions about whether to push or conserve, overtake or draft behind.
F1 has so much data for practically the same probabilistic framework, and I am
currently adapting this model for Solar Car and porting it to Rust so that it can take in
live telemetry data during a race and return recommendations to the driver in real-time this summer.
We placed 2nd last year; check back this summer, when we (hopefully) place first at the
Formula Sun Grand Prix!

\section{Data}

All data is from OpenF1 (\texttt{openf1.org}), using the 2026 Australian Grand
Prix (session key 11234, 56 laps, 20 race-finishing drivers).

\textbf{Variables:} car telemetry at 3.7\,Hz (speed km/h, throttle 0--100,
brake bool); lap records (number, duration, start timestamp); position snapshots
(driver, position, date); interval data (gap to leader per lap).

\textbf{Preprocessing:} Telemetry is fetched once per driver for the full
session ($\approx$31,000 points). Each lap is sliced by timestamp:
\[
  \text{tel}_{d,\ell} = \{p \in \text{tel}_d \mid t_\ell \le t_p \le t_\ell + \Delta_\ell + 1\}
\]
where $t_\ell$ is \texttt{date\_start} and $\Delta_\ell$ is lap duration. At
3.7\,Hz, a typical 87\,s lap yields $\approx$320 telemetry points. Position is
matched via nearest-prior-timestamp lookup since the endpoint returns
(date,\,position) pairs rather than (lap,\,position) pairs.

\section{Methods}

\subsection{Stage 1: Gaussian MLE for Corner Execution}

For each corner $c \in \{T1,\ldots,T16\}$ and data point
$f \in \{\text{speed, throttle, brake}\}$, we fit a Gaussian via MLE on the
fastest 25\% of laps across all drivers. This is because I assuming that fast laps
represent technically optimal execution. The MLE estimators are the closed-form
sample mean and variance:
\[
  \hat\mu_{c,f} = \frac{1}{n}\sum_i x_{c,f,i}, \qquad
  \hat\sigma^2_{c,f} = \frac{1}{n}\sum_i (x_{c,f,i} - \hat\mu_{c,f})^2
\]

A driver's corner score is the log-likelihood of their telemetry under the fitted Gaussian:
\[
LL(\theta) = \sum_{i=1}^{n} \left[
    -\log(\sqrt{2\pi\theta_1})
    - \frac{(x_i - \theta_0)^2}{2\theta_1}
\right]
\]
\begin{center}
    \textit{\small where $\theta_0 = \hat{\mu}$ (mean) and $\theta_1 = \hat{\sigma}^2$ (variance) are MLE parameters fit on the fastest 25\% of laps data.}
\end{center}

Higher score means execution closer to the optimal distribution. The lap score
sums all 16 corner scores. Corner importance is quantified by how much removing
each corner shifts the Naive Bayes score (Stage~2), averaged across all laps ---
this identifies the ``highest impact corner'' reported in the UI.

\subsection{Stage 2: Naive Bayes Lap Classifier}

We define $y \in \{\text{winner, non-winner}\}$ where ``winner'' labels the
fastest 15\% of laps across the full field. We fit Gaussian Naive Bayes via MLE
ffor $P(x_{c,d} \mid y)$ for each (corner, data point) pair. Under the conditional
independence assumption:
\[
  \log P(y \mid \mathbf{x}) = \log P(y) + \sum_{c,d} \log P(x_{c,d} \mid y)
\]
normalized using logs, as we learnt in lecture, to prevent underflow. The independence assumption
is clearly violated, since speed and throttle are correlated, but the
classifier still works well as a continuous ``winner-likeness'' score in $[0,1]$ that
I feed into Stage~3 as the corner execution evidence component.

\subsection{Stage 3: Bayesian Updating}

We maintain $P(\text{winner} = i)$ over all $n = 22$ drivers. After each lap $k$:

\begin{infobox}
\[
  P(\text{win}=i \mid \ell_{1:k}) \;\propto\; L_k(i)^T \cdot P(\text{win}=i \mid \ell_{1:k-1})
\]
\[
  L_k(i) = w_1 r_{\text{pos}} + w_2 r_{\text{gap}} + w_3 r_{\text{lap}} + w_4 r_{\text{corner}}
\]
\[
  T = 0.25 + 0.30\cdot\tfrac{k}{K}, \qquad \textstyle\sum_j w_j = 1
\]
\end{infobox}

where $r_{\cdot}$ denotes rank-normalization to $[0,1]$ across the field.
Weights shift over the race: $w_1$ (position) ramps $0.35 \to 0.60$; $w_4$
(corner) decays $0.25 \to 0.05$; $w_2, w_3 = 0.20$ constant. Three key design
choices require justification:

\begin{itemize}
\item \textbf{Rank normalization.} When I first combined the four signals, the position of the drivers, which is an integer between 1 and 22, was being drowned out by the corner log-likelihoods, which were large negative numbers. I went to AI for help and was presented with the idea rank normalization: each signal is sorted across all 22 drivers and mapped to $[0,1]$ based on relative standing, so the best driver on any metric gets 1.0 and the worst gets 0.0. This makes all four signals on the same scale before weighting.

\item \textbf{Temperature $T < 1$.} This is another technique I found through AI. Without it, 56 sequential Bayesian updates caused the posterior to collapse --- one driver would reach near 100\% probability within 10 laps and never recover, even if they then fell back. Raising the likelihood $L$ to a power $T < 1$ reduceds how hard each single lap can change the posterior. $T$ ramps upward as the race stabilizes and positions actually are decisive. Very cool technique!!

\item \textbf{Laplace smoothing.} As we learnt in class, a zero likelihood permanently zeroes the posterior, and it makes recovery impossible. $L$ is blended with $\varepsilon/n$ ($\varepsilon = 0.03$) before tempering, ensuring no driver's likelihood is ever exactly zero. This does come with a limitation, which I address in the Limitations section.
\end{itemize}

The prior is $90\%$ uniform $+ 10\%$ signal from 2025 average finishing
positions ($1/\text{pos}^{1.5}$, normalized), weak enough that observed data
dominates quickly.

\subsection{Counterfactual Inference}

This is a technique I discovered through research that was beyond the course material, but still very relevant to both CS109 content and this project. For each driver $d$ at each lap $k$, I wanted to compute $\Delta P$ for four interventions by running a hypothetical Bayesian update from the current posterior with one factor modified:

\begin{infobox}
\[
  \Delta P(d, c) = P(\text{win}=d \mid \tilde\ell_k) - P(\text{win}=d \mid \ell_k)
\]
Interventions $c$: gain 1 position; cut gap by 1.5\,s; match fastest lap
in field; match best performer at driver's personal weakest corner.
\end{infobox}

The four $\Delta P$ values are ranked for the best advice. Counterfactual causal inference doesn't give
me something that is correlated with winning the race, but gives me an actual action the driver can make
that increases the chance of winning  the most from  where we are right now. The
corner advice is the most specific: for every corner $c$, the system finds the
best performer, computes the improvement in overall corner score if driver
$d$ matched that benchmark, runs the hypothetical update, and reports the corner
with the highest $\Delta P$ as the targeted recommendation. This was the most interesting part of the entire project!

\section{Results}

Applied to the 2026 Australian Grand Prix (56 laps, 20 finishers), the system
produces the following at the final lap:

\begin{itemize}
  \item Russell (race winner, P1) reaches $\approx$65\% win probability by
    end of race; Antonelli (P2) $\approx$28\%; Leclerc (P3) $\approx$18\%.
  \item Russell's probability rises sharply after lap~10 as the position weight
    begins to dominate, consistent with him leading cleanly after the opening
    stint.
  \item The highest-impact corner (T8) is identified consistently, matching the
    sector where execution variance across the field is greatest.
  \item Counterfactual advice for Alex Albon at lap~20 correctly
    identifies gaining one position as the highest-leverage intervention, and he actually
    does do that in the next lap, which increases his win probabilty by 5\% . 
\end{itemize}

\section{Limitations}

Corner segmentation divides each lap's telemetry into 16 equal time windows
rather than GPS-aligned corners, which is an approximation. The Naive Bayes
conditional independence assumption is violated by correlated telemetry
data. Retired drivers (e.g.\ PIA, HUL) have posteriors frozen at
retirement due to Laplace smoothing, and it does not reflect the fact that they are retired. The nearest-prior-timestamp
position lookup occasionally misassigns position for laps with mid-lap
position changes.

\section{Solar Car Application}

The reason I built this for F1 is that the data is free and rich. But the real
goal is for Stanford's solar car: Azimuth (very cool name!). Every concept in this model maps directly ---
you just swap the telemetry variables:

\vspace{4pt}
\begin{tabular}{@{}ll@{}}
\toprule
\textbf{F1 concept} & \textbf{Solar car equivalent} \\
\midrule
Corner execution score & Road segment efficiency \\
Lap time               & Segment completion time \\
Gap to leader          & Time deficit to P1 \\
Throttle / brake       & Motor current, regen, braking \\
Fastest 25\% of laps   & Most efficient 25\% of segments \\
$\Delta P$ advice      & Push vs.\ conserve decision \\
\bottomrule
\end{tabular}
\vspace{4pt}

The one addition for solar is an \textit{energy state variable} --- the current
battery percentage. Unlike F1 where a driver can always push harder (or at least a lot more than our little solar car),
 Azimuth has a hard energy budget. If you spend too much early, you may not even finish. So
battery \% enters the prior as a constraint and updates the probability of
finishing the race at all. The counterfactual question becomes: should I push
harder here, or conserve energy? The model answers that in real time.

I am currently porting this to Rust so it can ingest the live telemetry that is being sent wirelessly to an InfluxDB setup. Right now, it gets 
data every half-second, and my goal is to make this model as fast as possible so that during a race, a strategy member can quickly
make decisions and return recommendations to the driver. We placed 2nd at the Formula Sun Grand Prix last year. The goal this summer is
first.
\section{Credits}

\textbf{Claude (Anthropic).} I could not have built this without Claude. I had
very little front-end experience going in, and Claude helped me build a majority of the
GUI. Beyond that, Claude helped me see ways to apply statistical principles I
already knew in ways I hadn't thought about --- most notably counterfactual
inference, which I would not have arrived at properly on my own. A lot of the
depth in this model came from my conversations with it, which helped me not only gain new knowledge, but also improve my understanding of the existing CS109 concepts I wanted to implement..

\textbf{Jeremy Hsieh.} Jeremy has been working alongside me on building the
telemetry pipeline for Stanford's solar car. As I port this model to Rust for
real-time use this summer, having access to real solar car telemetry data is
something we are figuring out together --- and his contributions there have been
invaluable.

\textbf{Juliette Woodrow.} Juliette suggested adding corner execution data to the
model. That one idea meaningfully improved accuracy and made the visualization
so much more interactive and interesting --- the entire corner-level breakdown in
the UI exists because of that conversation.

\section{Conclusion}

This project implements a probabilistic race intelligence system grounded
in CS109 concepts: Gaussian MLE, Naive Bayes lap classification, sequential
Bayesian updating, and the newly discovered but incredibly useful concept of
counterfactual inference for actionable strategic advice. The system runs on real
telemetry and produces win probability estimates consistent with the actual race
result. This summer, we race at the Formula Sun Grand Prix --- a chance to test
this model in the real world for the first time. But that is just the warmup.
In 2027, Stanford competes at the World Solar Challenge in Australia with a
brand new car. The probabilistic reasoning built here will be running in the
background, in real time, helping us win at the global level.

\end{document}
